% ZK-OTP: Action-Bound Zero-Knowledge Two-Factor Authorization for EOAs and Safes
% Draft for IACR ePrint / systems-security workshop submission.
\documentclass[11pt]{article}

\usepackage[margin=1in]{geometry}
\usepackage{amsmath,amssymb,amsthm}
\usepackage{booktabs}
\usepackage{listings}
\usepackage[hidelinks]{hyperref}
\usepackage{xcolor}

\theoremstyle{definition}
\newtheorem{definition}{Definition}
\theoremstyle{plain}
\newtheorem{theorem}{Theorem}
\newtheorem{lemma}{Lemma}
\newtheorem{claim}{Claim}

\newcommand{\Hash}{\mathsf{H}}       % Poseidon hash over BN254
\newcommand{\keccak}{\mathsf{keccak}}
\newcommand{\Fr}{\mathbb{F}_r}       % BN254 scalar field
\newcommand{\secret}{\mathsf{s}}
\newcommand{\salt}{\mathsf{r}}
\newcommand{\commit}{\mathsf{C}}
\newcommand{\chal}{\mathsf{ch}}
\newcommand{\nul}{\mathsf{N}}

\title{ZK-OTP: Action-Bound Zero-Knowledge Two-Factor Authorization\\
for EOAs (EIP-7702) and Smart-Contract Wallets}
\author{Rafael Escrich\thanks{Independent researcher. \texttt{rafael@unblockthechain.com}}}
\date{\today}

\begin{document}
\maketitle

\begin{abstract}
Second-factor authentication for blockchain accounts is today either off-chain
(custodial 2FA, defeated once the key is exported) or coarse (multisig $m$-of-$n$,
which authenticates \emph{signers}, not \emph{actions}). We present \textbf{ZK-OTP},
a zero-knowledge one-time-authorization scheme that proves knowledge of a
registered secret \emph{bound to a specific action}---the exact call batch, chain,
verifying contract, monotonic nonce, and expiry---without revealing the secret and
without an RFC-6238 shared-secret server. Our construction derives an
action-specific challenge in-circuit and a nullifier as a \emph{circuit output}
(not a prover input), eliminating a class of malleability and replay bugs. We show
how to attach ZK-OTP to a plain externally-owned account (EOA) via \textbf{EIP-7702}
delegation---adding a second factor \emph{without} changing or replacing the EOA
key---and to institutional custody via a \textbf{Safe Module}. The same
authentication circuit, with companion registration and rotation circuits, serves
both integrations. We
give a threat model, a security argument for action-binding, replay-, cross-chain-,
and malleability-resistance, and an implementation over Poseidon/BN254 with a
Groth16 proving path plus a migration path to transparent FRI-based proving.
Across the three circuits ($240$--$1074$ R1CS constraints), Groth16 proofs are
$256$~bytes, verify on-chain in $198$--$231$k gas, and generate in
$581$--$887$\,ms in the reproducible snarkjs pipeline.
\end{abstract}

\section{Introduction}
\label{sec:intro}

% --- Motivation ---
Account security on public blockchains is bimodal. Retail users hold a single
secp256k1 key whose theft is total and irreversible; there is no ``second factor''
that survives key export. Institutions use $m$-of-$n$ multisignature wallets (e.g.,
Safe), which raise the bar on \emph{who} signs but authenticate \emph{principals},
not \emph{operations}: a compromised threshold of signers authorizes anything.
Meanwhile, the 2FA that Web2 users expect---TOTP (RFC 6238), push, passkeys---is
either custodial (a server holds the shared secret) or not cryptographically bound
to the on-chain action being authorized.

% --- The opportunity: EIP-7702 ---
\textbf{EIP-7702} changes the design space. It lets an EOA set a persistent
delegation indicator pointing to contract code, so an
ordinary EOA can, for the first time, enforce arbitrary authorization logic
\emph{in its own storage context} while keeping its address and key. The
delegation remains in effect until replaced or cleared, making delegate
correctness, storage isolation, and recovery procedures part of the security
surface. This makes it
possible to add a real, on-chain second factor to the exact accounts that today
have none---without asking users to migrate to a smart-contract wallet.

% --- The gap ---
What has been missing is a second factor that is (i) \emph{action-bound}: the proof
authorizes one specific call batch, not a session or a signer; (ii) \emph{non-custodial}:
no server holds a shared secret; (iii) \emph{replay- and cross-chain-safe}; and
(iv) portable across a plain EOA (via 7702) and a Safe (via a Module), with one
authentication circuit family.

% --- Contributions ---
\paragraph{Contributions.}
\begin{enumerate}
  \item \textbf{An action-bound ZK-OTP scheme.} A zero-knowledge proof of knowledge
  of a registered secret whose challenge is the hash of
  $(\mathit{chainId}, \mathit{verifyingContract}, \mathit{actionHash}, \mathit{nonce},
  \mathit{deadline}, \commit)$, so a valid proof authorizes exactly one action on one
  chain before one deadline (\S\ref{sec:scheme}). The nullifier is derived
  \emph{in-circuit} and exposed as a public output, closing a malleability gap
  present in ``nullifier-as-input'' designs.
  \item \textbf{An EIP-7702 integration} (\texttt{ZkOtpDelegate}) that adds a ZK
  second factor to a plain EOA without replacing its key, with ERC-7201 namespaced
  storage and exhaustive public-signal pinning against replay (\S\ref{sec:onchain}).
  \item \textbf{A Safe Module variant} (\texttt{ZkOtpSafeModule}) that gates
  \texttt{execTransactionFromModule} on the same proof, giving institutions
  per-action ZK authorization on top of existing multisig (\S\ref{sec:safe}).
  \item \textbf{Registration and rotation circuits} that prevent commitment-griefing
  and enable key rotation without lockout (\S\ref{sec:circuits}).
  \item \textbf{A security analysis} against a defined threat model, and an
  implementation over Poseidon/BN254 with Groth16, evaluated on constraint count,
  proof size, on-chain verification gas, and end-to-end proof generation latency
  (\S\ref{sec:security}--\S\ref{sec:eval}). We discuss transparent FRI-based
  proving as the natural migration path for deployments that require no trusted
  setup or post-quantum assumptions.
\end{enumerate}

\section{Background}
\label{sec:background}

\paragraph{EIP-7702 delegation.} EIP-7702 introduces a set-code transaction whose
authorization tuple signs $(\mathit{chainId}, \mathit{address}, \mathit{nonce})$.
The protocol writes a delegation indicator into the EOA's code; subsequent calls
load code from the delegated address and execute it in the EOA's context. A
$\mathit{chainId}=0$ authorization is valid across chains, so ZK-OTP requires
wallets and relayers to use explicit non-zero chain identifiers.

\paragraph{ZK proof systems.} Groth16 gives succinct, constant-size proofs and
cheap on-chain verification, but it is pairing-based, quantum-vulnerable, and
requires a per-circuit trusted setup. Transparent FRI/STARK systems remove the
trusted setup and are the natural post-quantum migration path, at the cost of
larger proofs and higher verification cost.

\paragraph{Poseidon over BN254.} The implementation uses circomlib Poseidon~\cite{poseidon} over
BN254 (Table~\ref{tab:params}), matching the deployed Circom/Groth16 toolchain.
Changing the hash or field is a parameter migration: it requires regenerating
circuits, verifiers, test vectors, and benchmarks.

\paragraph{Safe modules and OTP.} Safe modules can execute transactions through
\texttt{execTransactionFromModule} once enabled by the Safe. Classical TOTP
assumes a shared secret and a time window; on-chain hash-chain OTP variants avoid
a server but consume or reveal a finite sequence. Neither form is naturally bound
to the full transaction semantics without additional protocol machinery.

\section{Threat Model and Goals}
\label{sec:threat}

\paragraph{System and assets.} A user controls an EOA key $k$ and a ZK secret
$\secret \in \Fr$. The protected asset is the authority to execute a specific call
batch from the EOA (7702) or Safe. The ZK factor is \emph{additive}: it does not
replace $k$.

\paragraph{Adversary.} We consider a PPT network adversary $\mathcal{A}$ who sees all
transactions and proofs (mempool), can reorder, replay, and submit arbitrary
transactions on any chain and contract, and may adaptively request honest
authentications for actions of its choice. Freshness below is always relative to
actions that were not queried to this honest-authentication oracle. We assume:
(i) the ZK secret $\secret$ is
generated uniformly in $\Fr$ and kept client-side (its compromise is a non-goal);
(ii) the Groth16 common reference string is generated by a ceremony with at least one
honest contributor (\S\ref{sec:prims}); a subverted CRS breaks knowledge-soundness
and motivates a transparent FRI migration path; (iii) the wallet faithfully sets
\texttt{actionHash} to the batch the user intends (action phishing is a UI non-goal);
and (iv) EIP-7702 delegation authorizations are pinned to a concrete
\texttt{chainId} (never the $\texttt{chainId}=0$ wildcard), enforced by the
relayer/wallet.

\paragraph{Security goals.}
\begin{enumerate}
  \item \textbf{Action-binding.} A valid proof authorizes exactly the call batch
  whose \texttt{actionHash} it commits to; it cannot be reused for a different batch.
  \item \textbf{Replay-resistance.} A proof is consumed by a monotonic nonce and a
  single-use nullifier; it cannot be replayed on the same chain.
  \item \textbf{Cross-chain replay-resistance.} The challenge binds \texttt{chainId}
  and \texttt{verifyingContract}; a proof valid on one chain/contract is invalid
  elsewhere.
  \item \textbf{Non-malleability.} The nullifier is uniquely determined by the
  witness (circuit output), so no third party can transform a valid proof into another
  accepted proof.
  \item \textbf{Knowledge-soundness of registration/rotation.} One cannot register
  or rotate to a commitment whose preimage one does not know.
  \item \textbf{Expiry.} A proof past its \texttt{deadline} is rejected.
\end{enumerate}

\paragraph{Non-goals (explicit).} Local device/secret compromise; action phishing
by a malicious frontend (a wallet-UI problem---the proof faithfully binds whatever
\texttt{actionHash} it is given); theft of the EOA key $k$ itself (an attacker with
$k$ can issue a new 7702 delegation to a non-ZK delegate---the factor is additive);
guardian/social recovery; and bridge/rollup-level attacks. See \S\ref{sec:limits}.

\section{The ZK-OTP Scheme}
\label{sec:scheme}

\subsection{Overview}
Three composable layers share the same authentication relation and verifier:
(1) a ZK-OTP circuit family; (2) \texttt{ZkOtpDelegate}, the EIP-7702 delegate; (3)
\texttt{ZkOtpSafeModule}. A user \emph{registers} a commitment
$\commit = \Hash(\secret, \salt)$ once, then \emph{authenticates} each action with a
fresh proof, and may \emph{rotate} the secret.

\subsection{Primitives}
\label{sec:prims}
We instantiate $\Hash$ as circomlib Poseidon over the BN254 scalar field $\Fr$ with the pinned
parameters of Table~\ref{tab:params}. The Groth16 path yields $\sim$256-byte proofs
verified in $\sim$250k gas. A future FRI/STARK instantiation would remove trusted
setup and improve post-quantum posture, at larger proof/verification cost. Plonky3~\cite{plonky3}
is one candidate implementation direction.

\begin{table}[t]
\centering
\caption{Pinned Poseidon/BN254 parameters.}
\label{tab:params}
\begin{tabular}{ll}
\toprule
Parameter & Value \\
\midrule
Field & BN254 $\Fr$ ($r = 2188\ldots495617$) \\
State width $t$ & 3 (2-input), 7 (6-input) \\
S-box $\alpha$ & 5 \\
Full rounds $R_F$ & 8 \\
Partial rounds $R_P$ & 56 ($t{=}3$), 60 ($t{=}7$) \\
\bottomrule
\end{tabular}
\end{table}

\subsection{Circuits}
\label{sec:circuits}

\paragraph{Authentication (\texttt{otp\_auth}).}
Public signals: $\mathit{chainId}, \mathit{verifyingContract}, \mathit{actionHash},
\mathit{nonce}, \mathit{deadline}, \commit, \nul$. Witness: $\secret, \salt$.
Constraints:
\begin{align}
  \commit &= \Hash(\secret, \salt), \label{eq:commit}\\
  \chal   &= \Hash(\mathit{chainId}, \mathit{verifyingContract}, \mathit{actionHash},
              \mathit{nonce}, \mathit{deadline}, \commit), \label{eq:chal}\\
  \nul    &= \Hash(\secret, \chal). \label{eq:null}
\end{align}
Crucially, $\nul$ is \emph{derived} by \eqref{eq:null} and exposed as a public
output; the contract reads it from the public signals rather than accepting it as an
input, so its value is uniquely fixed by the witness.

\paragraph{Registration (\texttt{otp\_register}).}
Public: $\commit, \mathit{bindingHash}$; witness: $\secret, \salt$. The circuit
constrains $\commit = \Hash(\secret, \salt)$. The contract separately pins
$\mathit{bindingHash}$ to
$\keccak(\commit, \mathit{userAddress}, \mathit{chainId}) \bmod |\Fr|$ before
calling the verifier. This forces a registrant to know the preimage of $\commit$
(preventing commitment-slot griefing), while the on-chain check binds the
registration calldata to an address and chain.

\paragraph{Rotation (\texttt{otp\_rotate}).}
Public signals: $\mathit{oldC}, \mathit{newC}, \mathit{chainId}, \mathit{verifyingContract},
\mathit{nonce}, \mathit{deadline}, \nul$; witness: old/new secret and salt.
Constraints prove knowledge of both preimages and derive the nullifier from the
\emph{old} secret, so rotation both proves current ownership and installs a
known new commitment---no lockout, no rotation to an unknown key.

\subsection{On-chain: the EIP-7702 Delegate}
\label{sec:onchain}
\texttt{ZkOtpDelegate} stores $(\commit, \mathit{nonce}, \text{used nullifiers})$ in
an ERC-7201 namespaced slot~\cite{erc7201} to avoid collisions under delegation. On
\texttt{executeWithProof(calls, proof, publicSignals)} it (i) recomputes each public
signal from runtime state (\texttt{block.chainid}, \texttt{address(this)},
$\keccak(\text{calls}) \bmod |\Fr|$, stored nonce, current time bound) and reverts on any
mismatch (\emph{public-signal pinning}); (ii) rejects a used nullifier; (iii) calls
the Groth16 verifier; (iv) marks the nullifier used and increments the nonce
\emph{before} executing \texttt{calls} (effects-before-interactions). Every
state-changing entry point is either \texttt{onlySelf} or proof-gated.

The call encoding is part of the trusted protocol surface: each batch item must
include the target address, ETH value, and calldata (and, for Safe, the Safe address
and module address). Binding calldata alone is insufficient, because value-bearing
calls would otherwise leave \texttt{msg.value} or per-call value outside the proved
statement. The resulting digest is reduced into $\Fr$ only after hashing a canonical
byte encoding; implementations should reject non-canonical encodings and document the
field-reduction rule used by both Solidity and Circom.

\subsection{Safe Module}
\label{sec:safe}
\texttt{ZkOtpSafeModule} gates \texttt{execTransactionFromModule} on the same
\texttt{otp\_auth} proof, so an institution's existing $m$-of-$n$ Safe additionally
requires a per-action ZK factor. The Safe setting has a different operational
surface from the EIP-7702 delegate: the Safe owner set, module enablement, and
module-disable procedures remain part of the institution's custody policy and
must be audited alongside the module.

\subsection{Off-chain prover}
\label{sec:prover}
The prover derives witness values from a client-held secret stored in an OS
keychain, secure enclave, HSM, or equivalent secret store. It computes the
Poseidon commitment/nullifier witness, generates the Groth16 proof, and passes
the proof plus public signals to the wallet or relayer. Before proof generation,
the wallet must decode the call batch into human-readable intent and display the
exact action whose hash will be bound into the circuit; otherwise ZK-OTP would
faithfully authorize a maliciously supplied \texttt{actionHash}.

\section{Security Analysis}
\label{sec:security}

\subsection{Assumptions}
\label{sec:assumptions}
Let $\Pi=(\mathsf{Setup},\mathsf{Prove},\mathsf{Vfy})$ be the Groth16 argument for an
NP relation $R$.

\begin{definition}[Knowledge soundness]
$\Pi$ is knowledge-sound if for every PPT prover $\mathcal{P}^*$ there is a PPT
extractor $\mathcal{E}$ such that, whenever $\mathcal{P}^*$ outputs $(x,\pi)$ with
$\mathsf{Vfy}(vk,x,\pi)=1$, $\mathcal{E}$ (given $\mathcal{P}^*$'s coins) outputs $w$
with $(x,w)\in R$, except with probability $\epsilon_{\mathrm{ks}}(\lambda)$.
\end{definition}

For Groth16, this is the standard knowledge-soundness assumption under the
algebraic/generic group style models used in succinct-argument analyses, conditioned
on a correctly generated CRS.

We model $\Hash$ (Poseidon over $\Fr$) as collision-resistant and preimage-resistant,
with advantages $\epsilon_{\mathrm{cr}},\epsilon_{\mathrm{pre}}$, and $\keccak$ as
collision-resistant. All named $\epsilon(\cdot)$ are negligible in $\lambda$ under a
correctly generated CRS (assumption (ii) of \S\ref{sec:threat}); a subverted CRS
voids $\epsilon_{\mathrm{ks}}$.

\subsection{Relations and on-chain acceptance}
The authentication circuit fixes the relation
\[
R_{\mathrm{auth}}=\Big\{\big(x=(\mathit{cid},\mathit{vc},\mathit{ah},\mathit{nc},
\mathit{dl},\commit,\nul);\, w=(\secret,\salt)\big):
\commit=\Hash(\secret,\salt)\,\wedge\,\nul=\Hash(\secret,\chal)\Big\},
\]
with $\chal=\Hash(\mathit{cid},\mathit{vc},\mathit{ah},\mathit{nc},\mathit{dl},\commit)$
as in \eqref{eq:chal}. $R_{\mathrm{reg}}$ fixes $\commit=\Hash(\secret,\salt)$;
$R_{\mathrm{rot}}$ fixes both preimages and a nullifier under the \emph{old} secret.

The delegate maintains state $(\mathit{cid}^*,\mathit{vc}^*,\mathit{nc}^*,\commit^*,
\mathsf{Used})$ and, at wall-clock $t$, \textbf{accepts} a submission
$(\mathsf{calls},x,\pi)$ iff
\begin{enumerate}\itemsep2pt
\item[\textbf{(P)}] \emph{pinning}: $\mathit{cid}=\mathit{cid}^*$, $\mathit{vc}=\mathit{vc}^*$,
$\mathit{ah}=\keccak(\mathsf{encodeCalls}(\mathsf{calls})) \bmod |\Fr|$, $\mathit{nc}=\mathit{nc}^*$,
$\commit=\commit^*$, and $t\le \mathit{dl}$;
\item[\textbf{(U)}] \emph{freshness}: $\nul\notin\mathsf{Used}$;
\item[\textbf{(V)}] \emph{validity}: $\mathsf{Vfy}(vk_{\mathrm{auth}},x,\pi)=1$.
\end{enumerate}
On acceptance it sets $\mathsf{Used}\mathrel{+}=\{\nul\}$, increments
$\mathit{nc}^*$, then executes $\mathsf{calls}$ (effects before interactions).
The monotonic nonce provides strict per-account ordering and state synchronization;
the nullifier provides proof-level replay protection independent of proof bytes and
remains necessary for rotation and for malleability-safe consumption semantics. A
future unordered variant could replace the monotonic nonce with independent action
identifiers or a nullifier bitmap, at the cost of a different concurrency model.

\subsection{Guarantees}

\begin{theorem}[Action unforgeability]
\label{thm:unforge}
Fix an honestly registered $\commit^*=\Hash(\secret,\salt)$ with $\secret$ hidden from
$\mathcal{A}$. After adaptively obtaining accepting submissions from an honest
prover for actions $\mathsf{calls}_1,\dots,\mathsf{calls}_q$, no PPT $\mathcal{A}$ causes the delegate to
accept a submission for a \emph{fresh} action $\mathsf{calls}^*$ (with
$\keccak(\mathsf{encodeCalls}(\mathsf{calls}^*))\neq
\keccak(\mathsf{encodeCalls}(\mathsf{calls}_i))\,\forall i$) except with
probability at most $\epsilon_{\mathrm{ks}}+\epsilon_{\mathrm{pre}}+\epsilon_{\mathrm{cr}}$.
\end{theorem}
\begin{proof}
A winning submission satisfies (P), so its statement has $\commit=\commit^*$ and
$\mathit{ah}=\keccak(\mathsf{encodeCalls}(\mathsf{calls}^*)) \bmod |\Fr|$, and satisfies (V). By knowledge soundness,
$\mathcal{E}$ extracts $(\secret',\salt')$ with $\commit^*=\Hash(\secret',\salt')$,
except with $\epsilon_{\mathrm{ks}}$. Two cases. If $(\secret',\salt')=(\secret,\salt)$,
then $\mathcal{A}$ (via $\mathcal{E}$) has produced $\secret$, a preimage of the given
$\commit^*$ it was never told---a preimage-finder with advantage bounded by
$\epsilon_{\mathrm{pre}}$. If $(\secret',\salt')\neq(\secret,\salt)$, the two pairs are
a collision on $\Hash$, bounded by $\epsilon_{\mathrm{cr}}$. Summing bounds the event.
Crucially, $\mathit{ah}$ is a \emph{pinned public signal}: the extracted witness
authorizes exactly $\mathsf{calls}^*$, so knowledge of $\secret$ is what a fresh
action costs.
\end{proof}

\begin{theorem}[Replay resistance]
\label{thm:replay}
No accepting submission can be accepted twice, and two accepting submissions for
distinct $(\mathit{ah},\mathit{nc})$ carry distinct nullifiers except with
probability $\epsilon_{\mathrm{cr}}$.
\end{theorem}
\begin{proof}
Acceptance adds $\nul$ to $\mathsf{Used}$; a resubmission of the same $x$ fails (U).
A submission reusing the proof under a new state fails (P) because $\mathit{nc}$ is
pinned to the incremented $\mathit{nc}^*$. Distinctness: $\nul=\Hash(\secret,\chal)$ and
$\chal$ binds $(\mathit{ah},\mathit{nc})$; equal nullifiers for distinct
$(\mathit{ah},\mathit{nc})$ yield an $\Hash$-collision, bounded by $\epsilon_{\mathrm{cr}}$,
so the single-use set never conflates two genuinely different authorizations.
\end{proof}

\begin{theorem}[Cross-chain / cross-contract resistance]
\label{thm:xchain}
A proof accepted at $(\mathit{cid}^*,\mathit{vc}^*)$ is rejected at any
$(\mathit{cid}',\mathit{vc}')\neq(\mathit{cid}^*,\mathit{vc}^*)$; forging a fresh
acceptance there reduces to Theorem~\ref{thm:unforge}.
\end{theorem}
\begin{proof}
$\mathit{cid},\mathit{vc}$ are pinned by (P), so replaying $(x,\pi)$ elsewhere fails
before verification. A \emph{new} accepting statement for $(\mathit{cid}',\mathit{vc}')$
has these values folded into $\chal$ and hence into $R_{\mathrm{auth}}$; producing it
without $\secret$ contradicts Theorem~\ref{thm:unforge}. Assumption~(iv)
(no $\mathit{chainId}=0$ delegation) prevents the orthogonal 7702 hazard of installing
the delegate cross-chain.
\end{proof}

\begin{theorem}[Non-malleability of authorization]
\label{thm:nonmal}
Groth16 proof re-randomization does not enable replay or nullifier substitution.
\end{theorem}
\begin{proof}
Freshness (U) and the used-set key on $\nul$, a \emph{public signal} in $x$, not on
the proof bytes $\pi$. Groth16 proofs are re-randomizable, so an adversary can derive
$\pi'\neq\pi$ for the same $x$; but $x$ (hence $\nul$) is unchanged, so (U) still
rejects. To obtain a \emph{different} accepted $\nul'$ for the same action, the
adversary needs a witness with $\Hash(\secret,\chal)=\nul'\neq\Hash(\secret,\chal)$---
impossible for a fixed witness, and finding a second witness is an $\Hash$-collision
($\epsilon_{\mathrm{cr}}$). Thus binding replay protection to the circuit-output
nullifier, rather than to $\pi$, sidesteps Groth16 malleability by construction.
\end{proof}

\begin{theorem}[Registration/rotation soundness]
\label{thm:reg}
An accepting \texttt{register} (resp.\ \texttt{rotate}) implies the submitter knows a
preimage of the (new) commitment, except with $\epsilon_{\mathrm{ks}}$.
\end{theorem}
\begin{proof}
By knowledge soundness on $R_{\mathrm{reg}}$ ($R_{\mathrm{rot}}$), extract
$(\secret,\salt)$ with $\commit=\Hash(\secret,\salt)$ (and, for rotation, the old
preimage witnessed by the nullifier under $\secret_{\mathrm{old}}$ plus
$\commit_{\mathrm{new}}=\Hash(\secret_{\mathrm{new}},\salt_{\mathrm{new}})$). Hence one
cannot claim a commitment slot for a value one cannot later authenticate with
(anti-griefing), nor rotate to an unknown key (no lockout). The domain tag
$\keccak(\commit,\mathit{account},\mathit{cid})\bmod \Fr$ is checked on-chain against
the public \texttt{bindingHash}, binding the commitment to one account and chain.
\end{proof}

\paragraph{Expiry and secrecy.} Expiry is immediate from (P): $t\le\mathit{dl}$.
Zero-knowledge of $\Pi$ keeps $\secret$ hidden across the $q$ observed proofs, so the
$\epsilon_{\mathrm{pre}}$ reduction in Theorem~\ref{thm:unforge} is against an
adversary that learns nothing about $\secret$ from transcripts.

\paragraph{Scope of the argument.} These are game-based reductions to $\Pi$'s
knowledge soundness and $\Hash,\keccak$ collision/preimage resistance. We do not claim
a UC treatment of the delegate state machine; the argument models the on-chain
contract as the explicit acceptance predicate above.

\section{Implementation and Evaluation}
\label{sec:eval}
We implement the three circuits in Circom~2.1.6 over BN254 with circomlib's Poseidon,
compile to R1CS, and generate a Groth16 proving/verifying key per circuit via a
Powers-of-Tau ($2^{12}$) plus per-circuit Phase-2 setup. The snarkjs-exported
Solidity verifiers are wired into the \texttt{ZkOtpDelegate} and
\texttt{ZkOtpSafeModule} contracts and exercised end-to-end: proofs generated
off-chain verify both off-chain and on-chain. Table~\ref{tab:eval} reports measured
figures; the full pipeline (constraint counts, setup, proof generation, gas) is
scripted and reproducible.

\begin{table}[t]
\centering
\caption{Measured evaluation. Gas via Foundry; proving time is the reproducible snarkjs
pipeline measurement (witness generation $+$ Groth16 proof) on a desktop/laptop
development machine.}
\label{tab:eval}
\begin{tabular}{lcccc}
\toprule
Circuit & R1CS constraints & Proof size & On-chain verify gas & Prove time \\
\midrule
\texttt{otp\_auth}     & 834  & 256\,B & 231{,}080 & 581\,ms \\
\texttt{otp\_register} & 240  & 256\,B & 197{,}504 & 683\,ms \\
\texttt{otp\_rotate}   & 1074 & 256\,B & 231{,}080 & 887\,ms \\
\bottomrule
\end{tabular}
\end{table}

All circuits are well under the $5{,}000$-constraint budget; Groth16 proofs are the
canonical eight field elements ($256$~bytes of calldata) and verify on-chain in
$\approx\!200$--$231$k gas (\texttt{register} is cheaper: two public signals vs.\
seven). The additional on-chain hashing for canonical \texttt{encodeCalls} is linear
in the encoded batch size and small relative to pairing verification for standard
interactive transactions. End-to-end proof generation in the reproducible snarkjs
pipeline takes
$581$--$887$\,ms. These figures are sufficient for interactive authorization, but
we leave mobile-device latency, transparent-FRI proof size/gas, and full
\texttt{executeWithProof} gas (verifier plus public-signal pinning, storage, and
scenario calls) to future evaluation.

\paragraph{Reproducibility and honest scope.} The measurements above use real
Groth16 proofs and exported Solidity verifiers; the end-to-end contract tests accept
valid proofs on-chain. The Phase-2 setup used here is a fixed-entropy
\emph{measurement} build; a production deployment requires a multi-contributor
ceremony (\S\ref{sec:prims}) or a transparent proving backend. This paper reports
a measured research prototype, not an audited production
system.

\section{Discussion and Limitations}
\label{sec:limits}
We restate, and where possible bound, what ZK-OTP does \emph{not} solve.

\paragraph{Device and secret compromise.} The security argument assumes $\secret$ is
secret (\S\ref{sec:threat}). Malware that reads $\secret$ from client storage defeats
the factor; zero knowledge protects the secret in transit and on-chain, not at rest.
Hardware-bound storage (Secure Enclave, YubiKey-class, or a Pico-HSM bridge)
\emph{mitigates} by making $\secret$ non-exportable and moving proving into the
device, but does not eliminate a fully compromised host.

\paragraph{Action phishing.} The proof faithfully binds whatever \texttt{actionHash}
it is given; if a malicious frontend asks the user to authorize batch $X$ but proves
batch $Y$, the on-chain check still passes for $Y$. Defense is wallet-side: the signer
must decode \texttt{actionHash} into a human-readable operation before proving. This is
a UI obligation we make explicit, not a contract property.

\paragraph{EOA key theft (the additivity boundary).} ZK-OTP is an \emph{additive}
factor. An attacker holding the EOA key $k$ can issue a fresh EIP-7702 authorization
delegating to a \emph{different}, non-ZK contract, bypassing the factor entirely. Thus
ZK-OTP raises the bar (two independent secrets now required for a protected action) but
does not, alone, survive full key theft; pairing it with a Safe whose owner set is not
the single EOA (\S\ref{sec:safe}) closes this for the institutional tier.

\paragraph{Recovery.} We deliberately omit guardian/social recovery: it would become
the weakest link and expand the threat surface. Key rotation (\S\ref{sec:circuits})
covers planned secret changes. In the Safe module, loss of $\secret$ without a prior
rotation is unrecoverable unless the Safe's owners disable or replace the module
through their normal governance path. In the EIP-7702 EOA setting, the base EOA key
$k$ is an ultimate recovery hatch: a legitimate user who still controls $k$ can clear
or replace the delegate, just as an attacker with $k$ can bypass it. This is the
same additivity boundary as above, not an extra recovery protocol.

\paragraph{Mempool exposure and deadlines.} A valid proof in the mempool is exposed
exactly as any signed transaction is; the single-use nullifier and short
\texttt{deadline} (we recommend $5$\,min for routine actions, $60$\,s for high-value)
bound the window. Front-running the \emph{action} is prevented by action-binding
(Theorem~\ref{thm:unforge}); front-running the \emph{inclusion} is a general L1 issue.

\paragraph{Quantum and the gas/PQ trade-off.} The Groth16 back end is pairing-based
and quantum-vulnerable, and it requires a trusted CRS. A transparent FRI/STARK
instantiation would remove both concerns at higher proof size and verification
gas. We therefore treat Groth16 as the measured high-frequency path and transparent
proving as a migration target for treasury or institutional operations where
no-trusted-setup assumptions dominate cost.

\paragraph{Cross-domain scope.} ZK-OTP protects EOA- and Safe-level authorization on
one chain/contract. Cross-chain and rollup-specific attacks (poisoned RPCs, sequencer
or validator compromise, bridge failures) are out of scope and orthogonal. The
\texttt{chainId} binding prevents ordinary cross-chain replay, but it does not
distinguish two forked histories that intentionally or temporarily share the same
\texttt{chainId}; deployments that care about such forks should include an additional
domain separator, such as a deployment-specific fork identifier or finalized genesis
anchor, in the action hash.

\section{Related Work}
\label{sec:related}

\paragraph{Account abstraction as a platform.} ERC-4337~\cite{erc4337} routes
\emph{UserOperations} through an EntryPoint and lets a smart-contract account run
arbitrary validation logic, enabling plugin architectures (session keys, spending
limits, guardians). It does not, however, retrofit an \emph{existing} EOA: a user
must adopt a 4337 smart account. EIP-7702~\cite{eip7702} closes exactly this gap by
letting a plain EOA delegate its code in-place. ZK-OTP targets the 7702 setting
directly---adding a cryptographic second factor to the accounts that hold the bulk of
retail funds today---while remaining deployable as a 4337/Safe module. Session-key and
guardian modules~\cite{safe} are \emph{policy} layers (who may act, up to what limit);
they are complementary to, not a substitute for, a per-action secret-knowledge factor.
For 4337-style deployment, ERC-7562 validation-scope rules~\cite{erc7562} are relevant:
ZK-OTP's validation is static and non-interactive from the bundler's perspective, but
an implementation must keep proof verification and storage reads within the allowed
validation surface and storage-access rules.

\paragraph{Passkeys and WebAuthn on-chain.} A growing line of AA wallets use
WebAuthn/FIDO2 passkeys~\cite{webauthn} as a signer, verifying secp256r1 signatures
on-chain via custom verifiers or the RIP-7212 precompile~\cite{rip7212}. Passkeys give
excellent hardware binding and UX, and a wallet can include transaction semantics in
the signed challenge, giving a public-key analogue of transaction confirmation. Their
primitive, however, is public-key/device authentication:
the account registers or verifies a public key and later checks signatures. ZK-OTP is
orthogonal and composable: a passkey can gate \emph{who} may prove, while the ZK-OTP
proof binds \emph{what} is authorized using a hidden secret and nullifier-based replay
protection. Hardware-bound secret storage (Secure Enclave, HSM) is our recommended
deployment for the ZK secret as well.

\paragraph{TOTP and hash-based OTP.} RFC-6238 TOTP~\cite{totp} is a shared-secret,
time-window authenticator: the verifier must know enough state to check the code, and
the code is not bound to a particular on-chain call. Hash-locked or Lamport-style
on-chain OTPs avoid a server, but they commit to a finite sequence and naturally
authorize possession of a token rather than a full transaction semantic. ZK-OTP keeps
no secret on-chain, is single-use via an in-circuit nullifier, is rotatable
(\S\ref{sec:circuits}), and is action-bound.

\paragraph{Nullifier-based ZK protocols.} Semaphore~\cite{semaphore} and mixer designs
such as Tornado Cash~\cite{tornado} popularised proving knowledge of a secret behind a
commitment while emitting a nullifier to prevent double-use, primarily for
\emph{anonymity} (set membership, unlinkable signaling/withdrawal). We reuse the
commitment/nullifier idiom but for a different end: the nullifier is derived from a
challenge over the \emph{action} (\S\ref{sec:scheme}), turning it into single-use
\emph{authorization} rather than anonymous signaling, and we add the register/rotate
lifecycle and the EOA/Safe integrations. Anonymity is explicitly \emph{not} a goal.
MACI~\cite{maci} is another important Ethereum ZK system using on-chain verification
to enforce correct private actions, but its goal is private, anti-collusion voting
under a coordinator model rather than per-account second-factor authorization.

\paragraph{ZK login and credential proofs.} zkLogin-style and keyless-account systems
prove possession of a Web2 (OAuth/OIDC/JWT) credential in zero knowledge to derive or
unlock an account~\cite{zklogin}. These target onboarding/identity from off-chain
credentials, not a per-action second factor over on-chain calls; the two can coexist
(ZK login to derive the account, ZK-OTP to authorize each sensitive action).

\paragraph{Threshold/MPC custody.} Multisig (Safe), threshold ECDSA, and commercial
MPC custody protect the \emph{key material} and distribute \emph{who} signs. They do
not by themselves prove knowledge of an independent second secret bound to the
operation semantics. ZK-OTP composes on top via a Safe module (\S\ref{sec:safe}),
adding an additive, per-action ZK factor over any such account.

\paragraph{Proof malleability.} Groth16~\cite{groth16} proofs are re-randomizable;
simulation-extractable SNARKs (e.g., GM17~\cite{gm17}) remove this at some cost. We
instead neutralise malleability at the protocol layer by keying replay protection on
the nullifier public signal rather than the proof bytes (Theorem~\ref{thm:nonmal}), so
the construction is not dependent on proof-byte uniqueness.

\paragraph{The gap.} To our knowledge, no prior system provides a second factor that
is simultaneously (i) \emph{action-bound} to a specific call batch/chain/contract/
nonce/deadline, (ii) \emph{non-custodial and secret-hiding} via zero knowledge,
(iii) attachable to a \emph{plain EOA} through EIP-7702 without migration or key
replacement, and (iv) realised by the same authentication circuit family across EOA
and institutional Safe deployments. ZK-OTP occupies this intersection.

\section{Conclusion}
ZK-OTP is a second factor for blockchain accounts that is bound to a specific action
rather than to a session or a signer, hides its secret in zero knowledge, and---via
EIP-7702---attaches to the plain EOAs that today have no second factor at all, using
the same authentication circuit family for an institutional Safe module. We gave game-based
security reductions to Groth16 knowledge-soundness and Poseidon/keccak
collision/preimage resistance, including a construction-level defense that makes
Groth16's proof malleability irrelevant, and a measured implementation:
$240$--$1074$ R1CS constraints, $256$-byte proofs, $198$--$231$k gas on-chain, and
desktop proof generation in $581$--$887$\,ms. The natural next steps are a production trusted-setup
ceremony, a transparent proving back end for the institutional tier, and an external
audit of the delegate and Safe module.

\begin{thebibliography}{99}
\bibitem{eip7702} V. Buterin, S. Wilson, A. Dietrichs, and lightclient. \emph{EIP-7702: Set Code for EOAs}. Ethereum Improvement Proposals, 2024. \url{https://eips.ethereum.org/EIPS/eip-7702}.
\bibitem{erc4337} V. Buterin et al. \emph{ERC-4337: Account Abstraction Using Alt Mempool}. Ethereum Improvement Proposals, 2021. \url{https://eips.ethereum.org/EIPS/eip-4337}.
\bibitem{groth16} J. Groth. \emph{On the Size of Pairing-Based Non-Interactive Arguments}. EUROCRYPT 2016.
\bibitem{poseidon} L. Grassi et al. \emph{Poseidon: A New Hash Function for Zero-Knowledge Proof Systems}. USENIX Security 2021.
\bibitem{semaphore} Semaphore. \emph{Semaphore: Anonymous signaling protocol}. \url{https://semaphore.pse.dev/}.
\bibitem{safe} Safe. \emph{Safe Smart Accounts and Modules}. \url{https://docs.safe.global/}.
\bibitem{erc7201} J. Amadi et al. \emph{ERC-7201: Namespaced Storage Layout}. Ethereum Improvement Proposals, 2023. \url{https://eips.ethereum.org/EIPS/eip-7201}.
\bibitem{erc7562} ERC-7562 authors. \emph{ERC-7562: Account Abstraction Validation Scope Rules}. Ethereum Improvement Proposals. \url{https://eips.ethereum.org/EIPS/eip-7562}.
\bibitem{plonky3} Polygon Zero. \emph{Plonky3}. \url{https://github.com/Plonky3/Plonky3}.
\bibitem{totp} D. M'Raihi et al. \emph{RFC 6238: TOTP}. IETF, 2011.
\bibitem{webauthn} W3C. \emph{Web Authentication: An API for accessing Public Key Credentials Level 2}. W3C Recommendation, 2021. \url{https://www.w3.org/TR/webauthn-2/}.
\bibitem{rip7212} A. Beregszaszi et al. \emph{RIP-7212: Precompile for secp256r1 curve support}. Ethereum Rollup Improvement Proposals. \url{https://github.com/ethereum/RIPs/blob/master/RIPS/rip-7212.md}.
\bibitem{tornado} A. Pertsev, R. Semenov, and R. Storm. \emph{Tornado Cash}. \url{https://tornado.cash/}.
\bibitem{maci} Privacy and Scaling Explorations. \emph{MACI: Minimum Anti-Collusion Infrastructure}. \url{https://maci.pse.dev/}.
\bibitem{zklogin} Mysten Labs. \emph{zkLogin}. \url{https://docs.sui.io/concepts/cryptography/zklogin}.
\bibitem{gm17} J. Groth, M. Maller. \emph{Snarky Signatures: Simulation-Extractable SNARKs}. CRYPTO 2017.
\end{thebibliography}

\end{document}
